\documentclass[12pt]{extarticle}

\usepackage{tabularx}
\usepackage{amsmath}
\usepackage{extsizes}
\usepackage{graphicx}
\usepackage{amssymb}
\usepackage[english,russian]{babel}
\usepackage[utf8]{inputenc}
\usepackage[a4paper]{geometry}
\geometry{top=20mm,bottom=30mm,left=20mm,right=15mm}
\usepackage[T1,T2A]{fontenc}
\usepackage{setspace}
\usepackage{titlesec}
\usepackage{indentfirst}
\usepackage{caption}
\renewcommand{\baselinestretch}{1.5}
\usepackage{spverbatim}
\usepackage{listings}
\usepackage{xcolor}
\usepackage{tabularray}
\lstset{extendedchars=\true}
\usepackage{hyperref}

\hypersetup{
    colorlinks=true,
    linkcolor=black,
    urlcolor=blue,
    pdftitle={Работа со справочниками городов},
}

\definecolor{sessionbg}{gray}{0.96}
\definecolor{sessionframe}{gray}{0.55}

\lstdefinestyle{session}{
    basicstyle=\ttfamily\small,
    breaklines=true,
    columns=fullflexible,
    keepspaces=true,
    frame=single,
    framesep=6pt,
    framerule=0.4pt,
    rulecolor=\color{sessionframe},
    backgroundcolor=\color{sessionbg},
    showstringspaces=false,
    upquote=true,
    aboveskip=8pt,
    belowskip=8pt,
    captionpos=b,
    xleftmargin=0pt,
    xrightmargin=0pt,
}

\newcommand{\cmd}[1]{\texttt{#1}}

\definecolor{codebg}{gray}{0.96}
\definecolor{comment}{gray}{0.45}
\definecolor{keyword}{RGB}{0, 75, 150}
\definecolor{string}{RGB}{150, 80, 30}
\definecolor{number}{gray}{0.5}

\lstdefinestyle{coolprint}{
  basicstyle   = \ttfamily\footnotesize,
  keywordstyle = \bfseries\color{keyword},
  commentstyle = \itshape\color{comment},
  stringstyle  = \color{string},
  backgroundcolor = \color{codebg},
  numbers      = left,
  numberstyle  = \footnotesize\color{number},
  numbersep    = 5pt,
  frame        = single,
  framerule    = 0.8pt,
  framesep     = 3pt,
  breaklines   = true,
  breakatwhitespace = true,
  tabsize      = 4,
  showspaces   = false,
  showstringspaces = false,
  xleftmargin  = 10pt,
  xrightmargin = 5pt,
  aboveskip    = \medskipamount,
  belowskip    = \medskipamount,
  captionpos   = b,
}

\usepackage{ragged2e}
\justifying

\titleformat{\section}{\centering\bfseries}{\thesection}{1em}{}
\titleformat{\subsection}{\centering\bfseries}{\thesubsection}{1em}{}

\begin{document}

\begin{center}
  \textbf{МИНОБРНАУКИ РОССИИ} \\
  \textbf{Санкт-Петербургский государственный} \\
  \textbf{электротехнический университет} \\
  \textbf{«ЛЭТИ» им. В.И. Ульянова (Ленина)} \\
  \textbf{Кафедра САПР} \\
\end{center}

\vspace{0.25\textheight}

\begin{center}
\textbf{Практическая работа №3} \\
\textbf{по дисциплине} \\
\textbf{«Объектно-ориентированное программирование»} \\
\end{center}

\vspace*{\fill}

\begin{center}
\begin{tabularx}{\textwidth}{ l X r }
 Студент гр. 5352 & \rule{8cm}{0.2mm} & Соболевский М.А. \\
 Преподаватель & \rule{8cm}{0.2mm} & Кулагин М.В. \\
\end{tabularx}
\end{center}

\vspace{3em}

\begin{center}
Санкт-Петербург\\
\the\year
\end{center}

\thispagestyle{empty}

\pagebreak

\section{Задание}

Разработать консольное приложение для работы с двумя
файлами-справочниками городов (в форматах CSV и XML). Приложение
принимает путь к файлу, разбирает его и выводит статистику:
список дублирующихся записей с количеством повторений,
распределение зданий по этажам в каждом городе, время обработки.
После вывода статистики приложение снова ждёт ввода пути или
команды завершения.

Парсеры CSV и XML должны иметь общий интерфейс, чтобы добавление
нового формата не требовало переписывать остальную логику.
Программа не должна падать на некорректном вводе; выход только по
команде пользователя.

\section{Спецификация программы}

\subsection{Общая архитектура}

Архитектура построена вокруг двух абстракций: \cmd{AbstractParser}
и \cmd{ParserSelector}. Диаграмма классов приведена на
рисунке~\ref{fig:classes}.

\begin{figure}[h]
\centering
\includegraphics[width=\textwidth]{build/report/assets/parser-classes}
\caption{Диаграмма классов}
\label{fig:classes}
\end{figure}

\cmd{AbstractParser}~--- абстрактный класс с единственным
публичным методом \cmd{parseFile}, в котором реализован общий
алгоритм: открыть источник, в цикле читать записи, закрыть
источник, собрать статистику. Три из этих шагов зависят от
формата файла и объявлены абстрактными: \cmd{initParser},
\cmd{parseRecord}, \cmd{close}. Конкретные парсеры
(\cmd{CSVParser} и \cmd{XMLParser}) переопределяют только их.

За счёт этого подсчёт статистики написан один раз, а отличия
форматов изолированы в наследниках. Если понадобится добавить
третий формат, достаточно написать ещё одного наследника с тремя
методами, остальной код трогать не придётся.

\cmd{ParserSelector<T>}~--- интерфейс для выбора парсера по ключу
типа~\cmd{T}. В текущей реализации ключом является расширение файла
(\cmd{ExtensionParserSelector}). Обобщённость позволяет при
необходимости сделать селектор по MIME-типу или по содержимому
файла, не меняя интерфейс.

\cmd{ParserSelectorFactory} собирает \cmd{ExtensionParserSelector}
из списка зарегистрированных парсеров. Список задан статически, а
расширение каждого парсера хранится в публичном статическом поле
\cmd{extension}. Фабрика читает это поле через рефлексию и
регистрирует парсер в селекторе. Если добавить новый парсер, нужно
только дописать его в списокd, регистрация произойдёт
автоматически.

\subsection{Модель данных}

Запись описывается классом \cmd{Record} с четырьмя полями:
город, улица, номер дома, этаж. Все поля объявлены \cmd{final},
чтобы запись нельзя было изменить после создания. Это нужно,
потому что \cmd{Record} используется как ключ в \cmd{HashMap}
и как элемент в \cmd{HashSet}. Методы \cmd{equals} и
\cmd{hashCode} реализованы вручную~--- сравнение по всем четырём
полям.

\cmd{Report} хранит результат обработки одного файла: длительность
обработки, словарь дубликатов (запись~$\to$~сколько раз
встретилась) и статистику по этажам. Формат вывода (метод
\cmd{toString}) отделён от самого расчёта. Если понадобится
другой формат, меняется только \cmd{toString}.

\subsection{Иерархия парсеров}

\cmd{CSVParser} использует библиотеку \cmd{fastcsv}. При
инициализации открывает файл и пропускает строку заголовков, при
каждом вызове \cmd{parseRecord} читает одну строку и превращает
её в \cmd{Record}.

\cmd{XMLParser} использует StAX~--- потоковый парсер из
стандартной библиотеки Java. При инициализации открывает поток,
при каждом вызове \cmd{parseRecord} ищет очередной элемент
\cmd{item} и читает из него атрибуты. Использован именно
потоковый парсер, а не DOM, потому что DOM загружает в память
всё дерево, а для справочника это лишнее. При большом файле
разница будет заметной.

Оба парсера читают файл по частям, а не целиком.

\subsection{Сбор статистики}

\cmd{parseFile} делает три вещи. Инициализирует парсер. В цикле
читает записи: каждую пробует добавить в \cmd{HashSet}. Если
добавилось, то это новая запись, нужно увеличить счётчик для её
этажа в соответствующем городе. Если не добавилось, то это
дубликат, нужно увеличить счётчик в словаре дубликатов. По
завершении закрывает парсер в \cmd{finally}.

Время замеряется через \cmd{Instant.now()} до открытия файла и
после закрытия. Разница оборачивается в \cmd{Duration} и попадает
в \cmd{Report}.

\section{Описание интерфейса пользователя программы}

\lstset{style=session}

Программа работает в режиме диалога. После запуска выводит
приглашение:

\renewcommand{\baselinestretch}{1}
\begin{lstlisting}
Enter filename:
\end{lstlisting}
\renewcommand{\baselinestretch}{1.5}

Пользователь вводит путь к файлу. Если расширение распознано и
файл успешно открывается, программа обрабатывает его и печатает
статистику. После вывода статистики снова появляется приглашение.

Завершение работы по вводу команды \cmd{exit} (или \cmd{quit})
либо по комбинации \texttt{Ctrl+D} (сигнал EOF на стандартном
вводе). В обоих случаях программа корректно закрывает ресурсы и
завершается с кодом~0.

При любой ошибке (несуществующий путь, неизвестное расширение,
повреждённый файл) программа печатает сообщение об ошибке и
возвращается к приглашению. Завершения по ошибке не происходит.

\section{Примеры работы}

Строки, которые вводит пользователь, набраны после приглашения
\cmd{Enter filename:}; остальные~--- вывод программы.

\subsection{Обработка CSV-файла}

\renewcommand{\baselinestretch}{1}
\begin{lstlisting}
Enter filename: data.csv
Time processing: 0.023s
Duplicates:
    {city: Барнаул, street: Дальняя улица, house: 56, floor: 2} repeated 2 times
Floors:
    Барнаул:
        1: 3 buildings
        2: 7 buildings
        3: 2 buildings
        4: 1 building
        5: 5 buildings
    Братск:
        1: 2 buildings
        2: 4 buildings
        3: 1 building
        4: 0 buildings
        5: 3 buildings
    ...

Enter filename:
\end{lstlisting}
\renewcommand{\baselinestretch}{1.5}

\subsection{Обработка XML-файла}

\renewcommand{\baselinestretch}{1}
\begin{lstlisting}
Enter filename: data.xml
Time processing: 0.041s
Duplicates:
    {city: Азов, street: Просека, улица, house: 156, floor: 3} repeated 3 times
Floors:
    Азов:
        1: 0 buildings
        2: 1 building
        3: 5 buildings
        4: 0 buildings
        5: 2 buildings
    ...
\end{lstlisting}
\renewcommand{\baselinestretch}{1.5}

\subsection{Ошибка: файл не найден}

\renewcommand{\baselinestretch}{1}
\begin{lstlisting}
Enter filename: nonexistent.csv
error: file not found: nonexistent.csv
Enter filename:
\end{lstlisting}
\renewcommand{\baselinestretch}{1.5}

\subsection{Ошибка: неизвестное расширение}

\renewcommand{\baselinestretch}{1}
\begin{lstlisting}
Enter filename: data.json
error: no parser registered for extension .json
Enter filename:
\end{lstlisting}
\renewcommand{\baselinestretch}{1.5}

\subsection{Завершение работы}

\renewcommand{\baselinestretch}{1}
\begin{lstlisting}
Enter filename: exit
$
\end{lstlisting}
\renewcommand{\baselinestretch}{1.5}

\section{Текст программы}
\lstset{style=coolprint}

\renewcommand{\baselinestretch}{1}
\subsection*{Main.java}
\lstinputlisting[language=Java]{src/main/java/impl/Main.java}

\subsection*{AbstractParser.java}
\lstinputlisting[language=Java]{src/main/java/impl/AbstractParser.java}

\subsection*{ParserSelector.java}
\lstinputlisting[language=Java]{src/main/java/impl/ParserSelector.java}

\subsection*{ExtensionParserSelector.java}
\lstinputlisting[language=Java]{src/main/java/impl/parserSelectors/ExtensionParserSelector.java}

\subsection*{ParserSelectorFactory.java}
\lstinputlisting[language=Java]{src/main/java/impl/ParserSelectorFactory.java}

\subsection*{Record.java}
\lstinputlisting[language=Java]{src/main/java/impl/Record.java}

\subsection*{Report.java}
\lstinputlisting[language=Java]{src/main/java/impl/Report.java}

\subsection*{CSVParser.java}
\lstinputlisting[language=Java]{src/main/java/impl/parsers/CSVParser.java}

\subsection*{XMLParser.java}
\lstinputlisting[language=Java]{src/main/java/impl/parsers/XMLParser.java}
\renewcommand{\baselinestretch}{1.5}

\section{Выводы}

В ходе работы было написано консольное приложение на Java для
разбора справочников городов в двух форматах. Основная задача
была не в самом разборе файлов (это стандартная операция,
поддержана библиотеками), а в том, чтобы правильно организовать
код так, чтобы форматы не «протекали» в остальную программу.

За счёт выделения абстрактного класса с общим алгоритмом удалось
избежать дублирования логики: подсчёт статистики и замер времени
написаны один раз и работают одинаково для CSV и XML. Добавление
нового формата сводится к реализации трёх методов~--- никакие
другие файлы трогать не нужно. То же самое с селектором парсеров:
интерфейс обобщённый, так что при необходимости можно сделать
другую реализацию (например, по содержимому файла) без изменения
кода, который её использует.

Из технических решений, о которых стоит упомянуть:
использование рефлексии для чтения поля \cmd{extension}
избавляет от отдельного списка расширений; неизменяемая модель
записи с корректными \cmd{equals} и \cmd{hashCode} позволяет
использовать её как ключ в \cmd{HashMap}; потоковые парсеры
(fastcsv и StAX) читают файл по частям и не требуют загружать
его в память целиком.

Ограничения текущей версии: обработка ошибок ввода-вывода
реализована через \cmd{try/catch} на верхнем уровне, подробной
диагностики нет; список регистрируемых парсеров задан статически,
автоматического сканирования по всему classpath нет. Обе вещи при
необходимости легко расширяются.

\end{document}
